\documentclass[10pt,a4paper]{amsart}
\usepackage[margin=19mm]{geometry}
\usepackage{amsmath,amssymb,mathtools}
\usepackage[scr=rsfs,cal=euler]{mathalfa}
\usepackage{xfrac}
\usepackage[unicode,hidelinks,bookmarks=false]{hyperref}
\newcommand{\ie}{i.e.\ }
\newcommand{\cf}{cf.\ }
\newcommand{\resp}{respectively\ }
\newcommand{\Subsection}{\subsection}
\providecommand{\nobreakdash}{\nobreak\hbox{-}}
\setlength{\emergencystretch}{2em}
\multlinegap0pt
\usepackage{bm}
\usepackage{bbm}
\usepackage{dsfont}
\usepackage{xspace}
\usepackage{cancel}
\usepackage{mathtools}
\usepackage{amsthm}
\makeatletter
\@ifundefined{theorem}{\newtheorem{theorem}{Theorem}}{}
\@ifundefined{lemma}{\newtheorem{lemma}[theorem]{Lemma}}{}
\makeatother
\renewcommand{\le}{\leqslant}\renewcommand{\ge}{\geqslant}
\title[The Syntax Of Polytopal Projections: From Permutohedra To As]{The Syntax Of Polytopal Projections: From Permutohedra To Associahedra}
\author{Milo\v{s} Ad\v{z}i\'c, Filip D. Jevti\'c}
\date{Source: arXiv:2605.25229v1 (CC BY 4.0)}
\begin{document}
\maketitle

\noindent\emph{Source and licence.} The Syntax Of Polytopal Projections: From Permutohedra To Associahedra. Version arXiv:2605.25229v1 (CC BY 4.0), distributed under the Creative Commons Attribution 4.0 International licence, as recorded in the arXiv licence metadata for this exact version and shown on the abstract page https://arxiv.org/abs/2605.25229v1. Full rights notice: licenses/arxiv-2605.25229-CC-BY-4.0.txt.\par\smallskip

\noindent\textit{Editorial scope.} Counted unit: the Lemma whose source label is \texttt{lem:f-word-normalization} in the article (source0.tex lines 992--1008, source label \texttt{lem:f-word-normalization}), delivered together with the article's own complete original proof of it (source0.tex lines 1010--1141, one \texttt{proof} environment of 132 lines). No mathematics is written, repaired or re-derived: statement and proof are the article's own text line for line. Required context retained uncounted: none. Every definition, display and label of those lines is retained unchanged. Omitted from this package and not redistributed: 1--991, 1009--1009, 1142--1660. Nothing inside a retained excerpt is omitted or altered: every retained body is delivered exactly as the article's lines are written, so no omission receipt of any kind accompanies this package. Citations: the retained excerpts carry no citation command at all, so no bibliography entry of the article is transcribed and no reference is redistributed. Reference adaptation: none was needed and none was made; every reference inside the delivered unit resolves inside it, so no unresolved reference or citation is produced. Printed slips kept literal: no printed word, symbol, hypothesis or number of the counted statement or of its proof was altered. Layout and class: the article's own class is \texttt{article} with packages including amsmath, amsthm, amssymb, url, graphicx, geometry, tikz, subcaption, mathtools, xcolor; the wrapper is the lane's pinned \texttt{amsart} editorial wrapper at 10pt on A4, which restates the theorem-like environments the retained text needs and adds the article's own macro definitions that the retained text needs (\texttt{\textbackslash le}), with extra wrapper packages (bm, bbm, dsfont, xspace, cancel, mathtools). The delivered numbering is the unit's own. Assets: no retained span contains a figure environment, a graphics-inclusion command, a PDF-page inclusion, a table or a TikZ picture; no image file is copied into this package, the package declares no asset dependency and no image file is redistributed with it. Licence: version arXiv:2605.25229v1 of this article is distributed under the Creative Commons Attribution 4.0 International licence, as recorded in the arXiv licence metadata for this exact version and shown on the abstract page https://arxiv.org/abs/2605.25229v1. A full rights notice is delivered at licenses/arxiv-2605.25229-CC-BY-4.0.txt. Characters: every retained excerpt is pure ASCII, so the delivered engine, XeLaTeX with its default Latin Modern text font, reports no missing character.
\begin{lemma}\label{lem:f-word-normalization}
Let \(a=a_1\cdots a_r\) be a nonempty word of distinct positive integers, and
let \(\kappa(a)=\kappa_1\cdots\kappa_r\) be the decreasing rearrangement of the
letters of \(a\) and write
\[
a=a_L\,\kappa_1\,a_R,
\]
where \(a_L\) and \(a_R\) are the subwords to the left and to the right of
\(\kappa_1\), respectively. 
Let $\mathbf B_L =_{\mathcal I} f(a_L)$ and $\mathbf B_R =_{\mathcal I} f(a_R)$. Then, $f(a)$ is equivalent modulo ${\mathcal I}$ to a term in one of the following forms:
\begin{itemize}
    \item [$\cdot$] if $a_L,a_R\neq\emptyset$, then $f(a)=_{\mathcal I}(\mathbf 2^{\kappa_1}\circ_2 \mathbf B_R)\circ_1 \mathbf B_L$;
    \item [$\cdot$] if \(a_L=\emptyset\), then $f(a)=_{\mathcal I}\mathbf 2^{\kappa_1}\circ_2 \mathbf B_R$;
    \item [$\cdot$] if \(a_R=\emptyset\), then $f(a)=_{\mathcal I}\mathbf 2^{\kappa_1}\circ_1 \mathbf B_L$;
    \item [$\cdot$] if \(a=\kappa_1\), then \(f(a)=\mathbf 2^{\kappa_1}\).
\end{itemize}
\end{lemma}

\begin{proof}
We argue by induction on \(r=\operatorname{len}(a)\). If \(r=1\), then \(a=\kappa_1\) and $f(a)=\mathbf 2^{\kappa_1}$, so the claim is immediate. Assume now that \(r>1\). Let \(a'\) be the word
obtained from \(a\) by deleting its smallest letter \(\kappa_r\). Then
\[
f(a)=f(a')\circ_{u_a(\kappa_r)+1}\mathbf 2^{\kappa_r}.
\]
The maximum of \(a'\) is still \(\kappa_1\). Write
\[
a'=a'_L\,\kappa_1\,a'_R.
\]
Since each axiom scheme preserves arity, congruent terms have the same arity.

We distinguish two cases.

\medskip
\noindent\emph{Case 1: \(\kappa_r\text{ occurs in } a_L\).} Then \(a'_R=a_R\), while \(a'_L\) is obtained from \(a_L\) by deleting
\(\kappa_r\). Moreover,
\[
u_a(\kappa_r)=u_{a_L}(\kappa_r),
\]
because every letter of \(a\) that is greater than \(\kappa_r\) and lies to its
left is already contained in \(a_L\).

\smallskip
\noindent If both \(a'_L\) and \(a'_R\) are nonempty, then by the inductive
hypothesis there exist terms \(\mathbf A_L,\mathbf A_R\in\mathcal L^I\) such that
\[
\mathbf A_L =_{\mathcal I} f(a'_L),\qquad \mathbf A_R =_{\mathcal I} f(a_R),
\]
and
\[
f(a') =_{\mathcal I} (\mathbf 2^{\kappa_1}\circ_2 \mathbf A_R)\circ_1 \mathbf A_L.
\]
Since \(\mathbf A_L =_{\mathcal I} f(a'_L)\), we have
\[
|\mathbf A_L|=|f(a'_L)|=\operatorname{len}(a'_L)+1=\operatorname{len}(a_L).
\]
Hence
\[
1\le u_{a_L}(\kappa_r)+1\le |\mathbf A_L|,
\]
so \emph{(assoc1)} applies:
\[
((\mathbf 2^{\kappa_1}\circ_2 \mathbf A_R)\circ_1 \mathbf A_L)\circ_{u_a(\kappa_r)+1}\mathbf 2^{\kappa_r}
=
(\mathbf 2^{\kappa_1}\circ_2 \mathbf A_R)\circ_1
\bigl(\mathbf A_L\circ_{u_{a_L}(\kappa_r)+1}\mathbf 2^{\kappa_r}\bigr).
\]
Set
\[
\mathbf B_L:=\mathbf A_L\circ_{u_{a_L}(\kappa_r)+1}\mathbf 2^{\kappa_r},
\qquad
\mathbf B_R:=\mathbf A_R.
\]
Since \(\mathbf A_L =_{\mathcal I} f(a'_L)\), it follows from the definition of \(f(a_L)\)
that
\[
\mathbf B_L =_{\mathcal I} f(a_L).
\]
Therefore
\[
f(a)=_{\mathcal I}(\mathbf 2^{\kappa_1}\circ_2 \mathbf B_R)\circ_1 \mathbf B_L.
\]

\smallskip
\noindent If either \(a'_L=\emptyset\) and \(a'_R\neq\emptyset\), or \(a'_R=\emptyset\) and \(a'_L\neq\emptyset\), we proceed analogously, and if \(a'=\kappa_1\) this case is trivial.

\medskip
\noindent\emph{Case 2: \(\kappa_r\text{ occurs in } a_R\).} Then \(a'_L=a_L\), while \(a'_R\) is obtained from \(a_R\) by deleting
\(\kappa_r\). Put
\[
s=u_{a_R}(\kappa_r)+1.
\]
Since every letter of \(a_L\), as well as \(\kappa_1\) itself, lies to the left
of \(\kappa_r\) and is greater than \(\kappa_r\), we have
\[
u_a(\kappa_r)=\operatorname{len}(a_L)+u_{a_R}(\kappa_r)+1.
\]
Again, we prove only the main case, since the other cases are similar.
\smallskip
\noindent If both \(a'_L\) and \(a'_R\) are nonempty, then by the inductive
hypothesis there exist terms \(\mathbf A_L,\mathbf A_R\in\mathcal L^I\) such that
\[
\mathbf A_L =_{\mathcal I} f(a_L),\qquad \mathbf A_R =_{\mathcal I} f(a'_R),
\]
and
\[
f(a') =_{\mathcal I} (\mathbf 2^{\kappa_1}\circ_2 \mathbf A_R)\circ_1 \mathbf A_L.
\]
Since \(\mathbf A_L =_{\mathcal I} f(a_L)\), we have
\[
|\mathbf A_L|=|f(a_L)|=\operatorname{len}(a_L)+1.
\]
Therefore
\[
u_a(\kappa_r)+1=|\mathbf A_L|+s.
\]
Applying \emph{(assoc2)} gives
\[
((\mathbf 2^{\kappa_1}\circ_2 \mathbf A_R)\circ_1 \mathbf A_L)\circ_{u_a(\kappa_r)+1}\mathbf 2^{\kappa_r}
=
\bigl((\mathbf 2^{\kappa_1}\circ_2 \mathbf A_R)\circ_{s+1}\mathbf 2^{\kappa_r}\bigr)\circ_1 \mathbf A_L.
\]
Now \(\mathbf A_R =_{\mathcal I} f(a'_R)\), so
\[
|\mathbf A_R|=|f(a'_R)|=\operatorname{len}(a'_R)+1=\operatorname{len}(a_R).
\]
Hence
\[
2\le s+1<2+|\mathbf A_R|,
\]
and \emph{(assoc1)} yields
\[
(\mathbf 2^{\kappa_1}\circ_2 \mathbf A_R)\circ_{s+1}\mathbf 2^{\kappa_r}
=
\mathbf 2^{\kappa_1}\circ_2\bigl(\mathbf A_R\circ_s \mathbf 2^{\kappa_r}\bigr).
\]
Set
\[
\mathbf B_L:=\mathbf A_L,
\qquad
\mathbf B_R:=\mathbf A_R\circ_s \mathbf 2^{\kappa_r}.
\]
Since \(\mathbf A_R =_{\mathcal I} f(a'_R)\), it follows from the definition of \(f(a_R)\) that
\[
\mathbf B_R =_{\mathcal I} f(a_R).
\]
Therefore
\[
f(a)=_{\mathcal I}(\mathbf 2^{\kappa_1}\circ_2 \mathbf B_R)\circ_1 \mathbf B_L.
\]
\end{proof}

\end{document}
